% !TEX TS-program = lualatex
Модуль edupra\_core реализует архитектуру системы обучения с подкреплением. Его структура представлена на диаграмме \refimage{edupra-core-struct} Архитектура разделена на несколько ключевых компонентов, каждый из которых инкапсулирует свою зону ответственности: модели (пакет models), агенты (пакет agents), утилиты (utils) и внешние зависимости, такие как библиотеки машинного обучения и игровое окружение.

\image[width=0.9\textwidth]{Строение модуля ядра машинного обучения}{edupra-core-struct}{edupra-core-struct}

В пакете models находятся классы, связанные с реализацией моделей обучения. Базовым является абстрактный класс BaseModel, который определяет основные параметры и методы, необходимые для работы с алгоритмами временных разниц. Он содержит методы инициализации, обновления весов, прямого прохода, сохранения и восстановления чекпоинтов, а также функции обучения и сравнения агентов. Класс BaseModel наследуется от keras.Model, что обеспечивает совместимость с экосистемой TensorFlow/Keras. Наследником BaseModel выступает TDGammon --- реализация нейросетевой модели, повторяющая архитектуру классического TD-Gammon. Она состоит из входного слоя (198 признаков), скрытого полносвязного слоя и одного выходного нейрона с функцией активации sigmoid.

Пакет agents реализует иерархию игровых агентов. Абстрактный класс Agent определяет общее поведение, включая бросок кубиков и выбор действия. Конкретные реализации включают RandomAgent (выбор случайного действия), HumanAgent (интерфейс для взаимодействия с человеком) и TDAgent (использует модель для принятия решений). Последние два класса используют экземпляры моделей, основанных на BaseModel, для оценки и выбора действий.

Пакет utils содержит вспомогательные компоненты, такие как модуль path с функцией ensure\_exists для работы с файловой системой, а также функцию evaluate\_agents, позволяющую массово тестировать агентов в игровом окружении и собирать статистику их эффективности.

Модуль также зависит от внешних библиотек. В первую очередь используется Keras --- часть TensorFlow, предоставляющая классы Model и Dense, необходимые для построения и обучения нейросетевых моделей. Кроме того, используется пользовательская реализация окружения BackgammonEnv, совместимая с API OpenAI Gym. Этот класс управляет взаимодействием между агентами и логикой игры, реализованной в классе Backgammon. Последний отвечает за генерацию допустимых ходов, выполнение действий, отслеживание состояния игры и определение победителя. Для описания состояния используется структура BackgammonState на основе namedtuple, что обеспечивает простоту сериализации и восстановления состояния игры. Класс Constants определяет ключевые параметры, такие как идентификаторы игроков, количество пунктов на доске, а также специальные зоны (BAR, OFF).
